\section{Background Knowledge}

\subsection{Markov Decision Process}
\begin{definition}[State, Action, Trajectory]
A sequence of actions and states with $T$ interactions with environment is defined as a trajectory $\tau = (s_0, a_0, s_1,a_1,s_2,a_2,\ldots, a_{T-1}, s_{T-1})$ 
\end{definition}

\begin{definition}[Reward and Return] 
% A single step reward $r_t$ is determined by the current state $s_t$ the action taken $a_t$ and the state $s_{t+1}$ after the action $a_{t}$ is taken. It is defined as $r_t := R(s_t,a_t,s_{t+1})$. 
At time step $t$, the agent observes state $s_t$, takes action $a_t$, and transitions to the next state $s_{t+1}$. The single-step reward is determined by the transition and is defined as
$$
r_t = R(s_t,a_t,s_{t+1}).
$$ 
The (discounted) \textbf{reward} on the trajectory $\tau$ is defined as 
$$
    R(\tau) = \sum_{t=0}^{t-1}\gamma^{t}r_t,
$$
where $\gamma\in(0,1]$ is the discount factor. The discounted \textbf{return} starting from time step $t$ is defined as
$$
    G_t(\tau) = \sum_{t'=t}^{T-1}\gamma^{t'-t}r_{t'}. 
$$
Equivalently, $G_t$ satisfies the recursive relation
$$
G_t = r_t + \gamma G_{t+1}.
$$
\end{definition}


\begin{definition}[State-Value Function and Value Function]
The value function for the policy $\pi$ is defined as
$$
% V^{\pi}(s_t) = \Embb_{\pi} [G_t | s_t] = \Embb_{a_t\sim \pi(.|s_t)}\Embb_{s_{t+1}\sim P(.|s_t, a_t)}\left[r_t + \gamma V^{\pi}(s_{t+1})\right].
\begin{aligned}
V^\pi(s_t)
&= \Embb_{\pi} [G_t \mid s_t]  = \Embb_{\pi} [r_t + \gamma G_{t+1} \mid s_t] = \mathbb{E}_{a_t\sim \pi(\cdot\mid s_t)}
\mathbb{E}_{s_{t+1}\sim P(\cdot\mid s_t,a_t)}
\left[
r_t + \gamma \Embb_{\pi}[G_{t+1}\mid s_{t+1}]
\right] \\
&= \mathbb{E}_{a_t\sim \pi(\cdot\mid s_t)}
\mathbb{E}_{s_{t+1}\sim P(\cdot\mid s_t,a_t)}
\left[
r_t + \gamma V^\pi(s_{t+1})
\right].
\end{aligned}
$$

The state-value function for the policy $\pi$ is defined as
$$
Q^{\pi}(s_t, a_t) = \Embb_{\pi} [G_t | s_t, a_t] = \Embb_{s_{t+1}\sim P(.|s_t, a_t)}\left[r_t + \gamma V^{\pi}(s_{t+1})\right].
$$
Therefore,
$$
V^{\pi}(s_t) = \Embb_{a_t\sim \pi(.|s_t)} [Q^{\pi}(s_t, a_t)]
$$
\end{definition}



\begin{definition}[Advantage]
The advantage function, by it’s definition, measures whether or not the action is better or worse than the policy’s default behavior.
    \begin{align}
        A^{\pi}(s_t, a_t) 
        & =  Q^{\pi}(s_t, a_t) - V^{\pi}(s_t)  = \Embb_{s_{t+1}\sim P(.|s_t, a_t)} \left[\underbrace{r_t + \gamma V^{\pi}(s_{t+1}) -  V^{\pi}(s_t)}_{\text{temporal difference}}\right]
    \end{align}
    
\end{definition}

\subsection{Useful Facts}

\begin{enumerate}
    \item  Probability of a trajectory 
        \begin{align}
        P(\tau | \theta)=\rho_0\left(s_0\right) \prod_{t=0}^{T-1} P\left(s_{t+1} | s_t, a_t\right) \pi_\theta\left(a_t | s_t\right) \label{eq:prob.of.a.trajectory}
        \end{align}
     \item  log derivative trick
         \begin{align}
        \grad_{\theta} P(\tau | \theta)=  P(\tau | \theta) \grad_\theta \log P(\tau | \theta)\label{eq:log.derivative.trick}
        \end{align}
    \item expected Grad-Log-Prob Lemma
        \begin{align}\label{eq:grad.log.prob}
            \Embb_{x\sim p(\cdot;\theta)} [\grad_{\theta} \log p(x;\theta)] = 0
        \end{align}

    
\end{enumerate}


\subsection{Glossaries}
\begin{definition}[Estimation of the $\operatorname{KL}$ divergence]\label{def:kl-est}
    There are different estimators of the KL divergence, which can be found \href{http://joschu.net/blog/kl-approx.html}{here}. Recall that $\operatorname{KL}(q||p)=E_{x\sim q}\left[\log \frac{q(x)}{p(x)}\right]$,
    
    \begin{itemize}
        \item $k_1$ estimator: $\frac{1}{n}\sum_{i=1}^{n} \log \frac{q(x_i)}{p(x_i)}$. Unbiased, high variance. Can be negative.
        \item $k_2$ estimator:  $\frac{1}{n}\sum_{i=1}^{n} \frac{1}{2}\left(\log \frac{p(x_i)}{q(x_i)}\right)^2$. biased, low variance. Always positive.
        \item  $k_3$ estimator: $\frac{p(x)}{q(x)} - 1 - \log \frac{p(x)}{q(x)}$. Always positive. Better than $k_2$? (non-negativity: $\log a \leq a - 1$ for all $a>0$).
    \end{itemize}
\end{definition}

\yd{To add: comparison with three estimators}